\chapter{Viscosity Solutions}

\section{Alexandroff Maximum Principle}

Fix a bounded connected domain $\Omega\subset\R^n$, $f\in C(\Omega)$,
and continuous coefficients satisfying
\[
\lambda|\xi|^2\leq a_{ij}(x)\xi_i\xi_j\leq\Lambda|\xi|^2.
\]
Write $Lu=a_{ij}D_{ij}u$.

\begin{definition}\label{hanlin:ch05-def-5-1}\dcref{Definition 5.1}\group{ch05-viscosity-solutions}
\source{han-lin-lecture-notes:page-0102}
A function $u\in C(\Omega)$ is a viscosity supersolution of $Lu=f$ if
every $\varphi\in C^2(\Omega)$ touching $u$ from below at
$x_0\in\Omega$ satisfies $L\varphi(x_0)\leq f(x_0)$. It is a viscosity
subsolution if every such test touching from above satisfies
$L\varphi(x_0)\geq f(x_0)$. A viscosity solution has both properties.
\end{definition}

\begin{remark}\label{hanlin:ch05-rem-5-2}\dcref{Remark 5.2}\group{ch05-viscosity-solutions}
\source{han-lin-lecture-notes:page-0102}
\uses{hanlin:ch05-def-5-1}
It is equivalent in Definition~\ref{hanlin:ch05-def-5-1} to use quadratic
test polynomials.
\end{remark}

\begin{remark}\label{hanlin:ch05-rem-5-3}\dcref{Remark 5.3}\group{ch05-viscosity-solutions}
\source{han-lin-lecture-notes:page-0102}
\uses{hanlin:ch05-def-5-1}
Classical supersolutions are viscosity supersolutions, and a $C^2$ viscosity
supersolution is classical. The analogous equivalence holds for subsolutions
and solutions.
\end{remark}
\begin{definition}\label{hanlin:ch05-def-5-4}\dcref{Definition 5.4}\group{ch05-viscosity-solutions}
\source{han-lin-lecture-notes:page-0102}
Fix $0<\lambda\leq\Lambda$ and $f\in C(\Omega)$. A function
$u\in C(\Omega)$ belongs to $S^+(\lambda,\Lambda,f)$ if every
$\varphi\in C^2(\Omega)$ touching it from below at $x_0$ satisfies
\[
\lambda\sum_{e_i>0}e_i+\Lambda\sum_{e_i<0}e_i\leq f(x_0),
\]
where the $e_i$ are the eigenvalues of $D^2\varphi(x_0)$. It belongs to
$S^-(\lambda,\Lambda,f)$ if every test touching from above satisfies
\[
\Lambda\sum_{e_i>0}e_i+\lambda\sum_{e_i<0}e_i\geq f(x_0).
\]
Set $S(\lambda,\Lambda,f)=S^+(\lambda,\Lambda,f)\cap
S^-(\lambda,\Lambda,f)$.
\end{definition}




\begin{remark}\label{hanlin:ch05-rem-5-5}\dcref{Remark 5.5}\group{ch05-viscosity-solutions}
\source{han-lin-lecture-notes:page-0103}
\uses{hanlin:ch05-def-5-4}
If
\[
a_{ij}D_{ij}u=f\quad\text{in }\Omega
\]
is uniformly elliptic with constants $\lambda,\Lambda$, then its viscosity
supersolutions lie in $S^+(\lambda,\Lambda,f)$. These extremal classes also
contain viscosity solutions of fully nonlinear uniformly elliptic equations.
\end{remark}

\begin{example}\label{hanlin:ch05-ex-5-6}\dcref{Example 5.6}\group{ch05-viscosity-solutions}
\source{han-lin-lecture-notes:page-0103}
\uses{hanlin:ch05-def-5-4}
Let $\mathcal A_{\lambda,\Lambda}$ be the symmetric matrices $A$ satisfying
\[
\lambda|\xi|^2\leq A_{ij}\xi_i\xi_j\leq\Lambda|\xi|^2.
\]
For symmetric $M$, define
\[
\mathcal M^-(M)=\inf_{A\in\mathcal A_{\lambda,\Lambda}}A_{ij}M_{ij},
\qquad
\mathcal M^+(M)=\sup_{A\in\mathcal A_{\lambda,\Lambda}}A_{ij}M_{ij}.
\]
If $e_1,\dots,e_n$ are the eigenvalues of $M$, then
\[
\mathcal M^-(M)=\lambda\sum_{e_i>0}e_i+\Lambda\sum_{e_i<0}e_i,
\qquad
\mathcal M^+(M)=\Lambda\sum_{e_i>0}e_i+\lambda\sum_{e_i<0}e_i.
\]
Thus $u\in S^+(\lambda,\Lambda,f)$ exactly when
\[
\mathcal M^-(D^2u)\leq f
\]
in the viscosity sense. For symmetric $M,N$,
\[
\mathcal M^-(M)+\mathcal M^-(N)
\leq\mathcal M^-(M+N)
\leq\mathcal M^+(M)+\mathcal M^-(N)
\leq\mathcal M^+(M+N)
\leq\mathcal M^+(M)+\mathcal M^+(N).
\]
\end{example}

For continuous $v$ on an open convex set $\Omega$, write
\[
\Gamma(v)(x)=\sup\{L(x):L\text{ affine and }L\leq v\text{ in }\Omega\}.
\]
Its contact set is $\{x\in\Omega:v(x)=\Gamma(v)(x)\}$.

\begin{lemma}\label{hanlin:ch05-lem-5-7}\dcref{Lemma 5.7}\group{ch05-viscosity-solutions}
\source{han-lin-lecture-notes:page-0104}
Let $u\in C^{1,1}(B_1)$ satisfy $u\geq0$ on $\partial B_1$, and let
$\Gamma_u$ be the convex envelope of $-u^-=\min\{u,0\}$. Then
\[
\sup_{B_1}u^-
\leq c(n)\left(\int_{B_1\cap\{u=\Gamma_u\}}\det D^2u\right)^{1/n}.
\]
\end{lemma}

\begin{theorem}\label{hanlin:ch05-thm-5-8}\dcref{Theorem 5.8}\group{ch05-viscosity-solutions}
\source{han-lin-lecture-notes:page-0104,han-lin-lecture-notes:page-0105}
\uses{hanlin:ch05-def-5-4, hanlin:ch05-lem-5-7}
Suppose $u\in S^+(\lambda,\Lambda,f)$ in $B_1$, $u\geq0$ on
$\partial B_1$, and $f\in C(\Omega)$. If $\Gamma_u$ is the convex envelope
of $-u^-$, then
\[
\sup_{B_1}u^-
\leq c\left(\int_{B_1\cap\{u=\Gamma_u\}}(f^+)^n\right)^{1/n},
\]
where $c$ depends only on $n,\lambda$, and $\Lambda$.
\end{theorem}

\begin{proof}
\sketch Viscosity touching tests control the Hessian of $\Gamma_u$ on the
contact set. Apply Lemma~\ref{hanlin:ch05-lem-5-7} to the resulting
determinant bound.
\end{proof}

\section{Harnack Inequality and \Calpha Bounds}

\begin{lemma}\label{hanlin:ch05-lem-5-9}\dcref{Lemma 5.9}\group{ch05-viscosity-solutions}
\source{han-lin-lecture-notes:page-0105,han-lin-lecture-notes:page-0106}
\uses{hanlin:weaksol-part1-cz}
Let $A\subset B\subset Q_1$ be measurable and let $\delta\in(0,1)$.
Assume
\begin{itemize}
\item[(i)] $|A|<\delta$;
\item[(ii)] if a dyadic cube $Q$ satisfies $|A\cap Q|\geq\delta|Q|$,
then its predecessor $\widetilde Q$ lies in $B$.
\end{itemize}
Then $|A|\leq\delta|B|$.
\end{lemma}

\begin{proof}
\sketch Apply Lemma~\ref{hanlin:weaksol-part1-cz} to $\chi_A$ and sum
the predecessor-cube estimates inside $B$.
\end{proof}



\begin{theorem}\label{hanlin:ch05-thm-5-10}\dcref{Theorem 5.10}\group{ch05-viscosity-solutions}
\source{han-lin-lecture-notes:page-0106,han-lin-lecture-notes:page-0107}
\uses{hanlin:ch05-lem-5-12}
If $u\in\Sclass(\lambda,\Lambda,f)$ in $\B_1$, $u\geq0$, and
$f\in C(\B_1)$, then
\[
\sup_{\B_{1/2}}u
\leq c\left(\inf_{\B_{1/2}}u+\|f\|_{L^n(\B_1)}\right),
\]
where $c$ depends only on $n,\lambda$, and $\Lambda$.
\end{theorem}
\begin{proof}
\sketch After scaling to $Q_{4\sqrt n}$, normalize by
\[
u_\delta=
\frac{u}{\inf_{Q_{1/4}}u+\delta+
\varepsilon_0^{-1}\|f\|_{L^n(Q_{4\sqrt n})}}.
\]
Lemma~\ref{hanlin:ch05-lem-5-12} applies to $u_\delta$. Let
$\delta\downarrow0$ and cover $B_{1/2}$ by scaled cubes.
\end{proof}

\begin{corollary}\label{hanlin:ch05-cor-5-11}\dcref{Corollary 5.11}\group{ch05-viscosity-solutions}
\source{han-lin-lecture-notes:page-0106}
\uses{hanlin:ch05-thm-5-10}
If $u\in\Sclass(\lambda,\Lambda,f)$ in $\B_1$ with $f\in C(\B_1)$,
then $u\in\Calpha(\B_1)$ for some $\alpha\in(0,1)$ depending only on
$n,\lambda$, and $\Lambda$, and
\[
|u(x)-u(y)|\leq c|x-y|^\alpha
\left(\sup_{\B_1}|u|+\|f\|_{L^n(\B_1)}\right)
\qquad(x,y\in\B_{1/2}),
\]
where $c$ has the same parameter dependence.
\end{corollary}
\begin{proof}
\sketch On each ball, apply Theorem~\ref{hanlin:ch05-thm-5-10} to the
nonnegative functions obtained by subtracting $u$ from its supremum and its
infimum from $u$. This contracts the oscillation up to the scaled $L^n$ norm
of $f$; iteration on concentric balls gives the stated exponent and estimate.
\end{proof}

\begin{lemma}\label{hanlin:ch05-lem-5-12}\dcref{Lemma 5.12}\group{ch05-viscosity-solutions}
\source{han-lin-lecture-notes:page-0106,han-lin-lecture-notes:page-0109,han-lin-lecture-notes:page-0110,han-lin-lecture-notes:page-0111}
\uses{hanlin:ch05-lem-5-13, hanlin:ch05-lem-5-14}
There are $\varepsilon_0,C>0$, depending only on $n,\lambda,\Lambda$,
such that the following holds. If
$u\in\Sclass(\lambda,\Lambda,f)$ in $Q_{4\sqrt n}$, $u\geq0$,
$f\in C(Q_{4\sqrt n})$, and
\[
\inf_{Q_{1/4}}u\leq1,
\qquad
\|f\|_{L^n(Q_{4\sqrt n})}\leq\varepsilon_0,
\]
then $\sup_{Q_{1/4}}u\leq C$.
\end{lemma}

\begin{proof}
\sketch Combine Lemma~\ref{hanlin:ch05-lem-5-13} with the distribution
decay in Lemma~\ref{hanlin:ch05-lem-5-14}. A point-selection iteration
would otherwise produce an unbounded sequence inside the controlled cube.
\end{proof}

\begin{lemma}\label{hanlin:ch05-lem-5-13}\dcref{Lemma 5.13}\group{ch05-viscosity-solutions}
\source{han-lin-lecture-notes:page-0107,han-lin-lecture-notes:page-0108}
\uses{hanlin:ch05-def-5-4, hanlin:ch05-thm-5-8}
There are $\varepsilon_0>0$, $\mu\in(0,1)$, and $M>1$, depending
only on $n,\lambda,\Lambda$, such that if
$u\in\Splus(\lambda,\Lambda,f)$ in $B_{2\sqrt n}$,
$f\in C(B_{2\sqrt n})$, and
\[
u\geq0,\qquad \inf_{Q_3}u\leq1,\qquad
\|f\|_{L^n(B_{2\sqrt n})}\leq\varepsilon_0,
\]
then $|\{u\leq M\}\cap Q_1|>\mu$.
\end{lemma}

\begin{proof}
\sketch Add a concave barrier and apply
Theorem~\ref{hanlin:ch05-thm-5-8} to its convex envelope. The contact set
is contained in a fixed sublevel set of $u$.
\end{proof}

\begin{lemma}\label{hanlin:ch05-lem-5-14}\dcref{Lemma 5.14}\group{ch05-viscosity-solutions}
\source{han-lin-lecture-notes:page-0108,han-lin-lecture-notes:page-0109}
\uses{hanlin:ch05-lem-5-9, hanlin:ch05-lem-5-13}
There are $\varepsilon_0,\varepsilon,C>0$, depending only on
$n,\lambda,\Lambda$, such that if
$u\in\Splus(\lambda,\Lambda,f)$ in $B_{2\sqrt n}$,
$f\in C(B_{2\sqrt n})$, and
\[
u\geq0,\qquad \inf_{Q_3}u\leq1,\qquad
\|f\|_{L^n(B_{2\sqrt n})}\leq\varepsilon_0,
\]
then, for every $t>0$,
\[
|\{u\geq t\}\cap Q_1|\leq Ct^{-\varepsilon}.
\]
\end{lemma}

\begin{proof}
\sketch Lemmas~\ref{hanlin:ch05-lem-5-9} and
\ref{hanlin:ch05-lem-5-13} give geometric decay of dyadic superlevel sets.
Reparameterizing the levels yields the power estimate.
\end{proof}



For bounded $\Omega$, continuous $u$, and $M>0$, let
$\Gm{M}{u,\Omega}$ be the points admitting a touching paraboloid from
below of opening $M$, and let $\Gp{M}{u,\Omega}$ be defined analogously
from above. Set
\[
\G{M}{u,\Omega}=\Gm{M}{u,\Omega}\cap\Gp{M}{u,\Omega},
\qquad
\Am{M}{u,\Omega}=\Omega\setminus\Gm{M}{u,\Omega},
\]
with analogous notation $\Ap{M}{u,\Omega}$ and
$\A{M}{u,\Omega}=\Omega\setminus\G{M}{u,\Omega}$.

\begin{lemma}\label{hanlin:ch05-lem-5-15}\dcref{Lemma 5.15}\group{ch05-viscosity-solutions}
\source{han-lin-lecture-notes:page-0111,han-lin-lecture-notes:page-0112,han-lin-lecture-notes:page-0113,han-lin-lecture-notes:page-0114}
\uses{hanlin:ch05-lem-5-9, hanlin:ch05-lem-5-16}
Let $B_{6\sqrt n}\subset\Omega$ with $\Omega$ bounded. There are
$\delta_0,\mu,C>0$, depending only on $n,\lambda,\Lambda$, such that if
$u\in S^+(\lambda,\Lambda,f)$ in $B_{6\sqrt n}$,
$f\in C(B_{6\sqrt n})$, $|u|\leq1$ in $\Omega$, and
$\|f\|_{L^n(B_{6\sqrt n})}\leq\delta_0$, then
\[
|\Am{t}{u,\Omega}\cap Q_1|\leq Ct^{-\mu}\qquad(t>0).
\]
If also $u\in S(\lambda,\Lambda,f)$ in $B_{6\sqrt n}$, then
\[
|\A{t}{u,\Omega}\cap Q_1|\leq Ct^{-\mu}\qquad(t>0).
\]
\end{lemma}

\begin{proof}
\sketch Lemma~\ref{hanlin:ch05-lem-5-16} supplies a uniform contact-set
gain. Iterate the complementary bad-set estimate over dyadic cubes and use
the same argument for $-u$ to obtain the two-sided conclusion.
\end{proof}

\begin{lemma}\label{hanlin:ch05-lem-5-16}\dcref{Lemma 5.16}\group{ch05-viscosity-solutions}
\source{han-lin-lecture-notes:page-0112}
\uses{hanlin:ch05-lem-5-13, hanlin:ch05-thm-5-8}
Let $B_{6\sqrt n}\subset\Omega$ with $\Omega$ bounded, and suppose
$u\in S^+(\lambda,\Lambda,f)$ in $B_{6\sqrt n}$ with
$f\in C(B_{6\sqrt n})$. There are $\sigma\in(0,1)$, $\delta_0>0$,
and $M>1$, depending only on $n,\lambda,\Lambda$, such that
\[
\|f\|_{L^n(B_{6\sqrt n})}\leq\delta_0,\qquad
\Gm{1}{u,\Omega}\cap Q_3\neq\emptyset
\]
implies $|\Gm{M}{u,\Omega}\cap Q_1|\geq1-\sigma$.
\end{lemma}

\begin{proof}
\sketch Combine the barrier construction of
Lemma~\ref{hanlin:ch05-lem-5-13} with
Theorem~\ref{hanlin:ch05-thm-5-8} to enlarge the lower contact set.
\end{proof}


\begin{remark}\label{hanlin:ch05-rem-5-17}\dcref{Remark 5.17}\group{ch05-viscosity-solutions}
\source{han-lin-lecture-notes:page-0114}
\uses{hanlin:ch05-lem-5-15}
For $u\in S(\lambda,\Lambda,f)$, polynomial decay of
\[
\mu(t)=|A_t(u,\Omega)\cap Q_1|
\]
implies $D^2u\in L^p(Q_1)$ for sufficiently small $p>0$, depending only
on $n,\lambda,\Lambda$. Larger exponents require faster decay.
\end{remark}

\section{Schauder Estimates}

Throughout this section, $a_{ij},f\in C(B_1)$ and
\[
\lambda|\xi|^2\leq a_{ij}(x)\xi_i\xi_j\leq\Lambda|\xi|^2.
\]

\begin{lemma}\label{hanlin:ch05-lem-5-18}\dcref{Lemma 5.18}\group{ch05-viscosity-solutions}
\source{han-lin-lecture-notes:page-0115}
\uses{hanlin:ch05-thm-5-8, hanlin:ch05-cor-5-11, hanlin:lemma-holder-boundary}
Let $u\in C(B_1)$ solve
\[
a_{ij}D_{ij}u=f\quad\text{in }B_1,
\]
in the viscosity sense, with $|u|\leq1$. If $0<\varepsilon<1/16$ and
\[
\|a_{ij}-a_{ij}(0)\|_{L^n(B_{3/4})}\leq\varepsilon,
\]
then some $h\in C(\overline{B_{3/4}})$ satisfies
\[
a_{ij}(0)D_{ij}h=0,\qquad |h|\leq1\quad\text{in }B_{3/4},
\]
and
\[
\|u-h\|_{L^\infty(B_{1/2})}\leq C\{\varepsilon^\gamma+\|f\|_{L^n(B_1)}\},
\]
where $C>0$ and $\gamma\in(0,1)$ depend only on $n,\lambda,\Lambda$.
\end{lemma}

\begin{proof}
\sketch Solve the frozen-coefficient Dirichlet problem with boundary data
$u$. Corollary~\ref{hanlin:ch05-cor-5-11} and
Lemma~\ref{hanlin:lemma-holder-boundary} control the boundary discrepancy on a
slightly smaller ball, while interior estimates control $D^2h$. Apply
Theorem~\ref{hanlin:ch05-thm-5-8} to $u-h$, whose right-hand side is
$f-(a_{ij}-a_{ij}(0))D_{ij}h$, and optimize the boundary-layer width.
\end{proof}


\begin{definition}\label{hanlin:ch05-def-5-19}\dcref{Definition 5.19}\group{ch05-viscosity-solutions}
\source{han-lin-lecture-notes:page-0116}
A function $g$ is $L^n$-Hölder continuous at $0$ with exponent $\alpha$ if
\[
[g]_{C^\alpha_{L^n}}(0)
:=\sup_{0<r<1}r^{-\alpha}
\left(\frac1{|B_r|}\int_{B_r}|g-g(0)|^n\right)^{1/n}<\infty.
\]
\end{definition}

\begin{theorem}\label{hanlin:ch05-thm-5-20}\dcref{Theorem 5.20}\group{ch05-viscosity-solutions}
\source{han-lin-lecture-notes:page-0116,han-lin-lecture-notes:page-0117,han-lin-lecture-notes:page-0118}
\uses{hanlin:ch05-def-5-19, hanlin:ch05-lem-5-18}
Let $u\in C(B_1)$ solve
\[
a_{ij}D_{ij}u=f\quad\text{in }B_1
\]
in the viscosity sense. If $a_{ij}$ and $f$ are $L^n$-Hölder at $0$ with
the same exponent $\alpha\in(0,1)$, then there is a quadratic polynomial
$P$ such that
\[
\|u-P\|_{L^\infty(B_r)}\leq Cr^{2+\alpha}\qquad(0<r<1),
\]
\[
|P(0)|+|DP(0)|+|D^2P(0)|\le C_*,
\]
\[
C_*\le C\{|u|_{L^\infty(B_1)}+|f(0)|+[f]_{C^\alpha_{L^n}}(0)\},
\]
where $C$ depends only on $n,\lambda,\Lambda,\alpha$ and
$[a_{ij}]_{C^\alpha_{L^n}}(0)$. In particular, $u$ is $C^{2,\alpha}$ at
$0$.
\end{theorem}

\begin{proof}
\sketch Subtract a quadratic that accounts for $f(0)$, then normalize the
solution and the coefficient oscillation. Iterating
Lemma~\ref{hanlin:ch05-lem-5-18} on geometrically shrinking balls produces
quadratics whose constant, linear, and Hessian coefficients form convergent
sequences. Their limit is $P$, and rescaling gives the $r^{2+\alpha}$ error
and coefficient bounds.
\end{proof}


\begin{theorem}\label{hanlin:ch05-thm-5-21}\dcref{Theorem 5.21}\group{ch05-viscosity-solutions}
\source{han-lin-lecture-notes:page-0118}
\uses{hanlin:ch05-thm-5-20}
Let $u\in C(B_1)$ be a viscosity solution of
\[
a_{ij}D_{ij}u=f\quad\text{in }B_1.
\]
For each $\alpha\in(0,1)$ there is $\theta>0$, depending only on
$n,\lambda,\Lambda,\alpha$, such that if
\[
\left(\frac1{|B_r|}\int_{B_r}|a_{ij}-a_{ij}(0)|^n\right)^{1/n}
\leq\theta\qquad(0<r\leq1),
\]
then an affine function $L$ satisfies
\[
\|u-L\|_{L^\infty(B_r)}\leq C_*r^{1+\alpha}\qquad(0<r<1),
\]
\[
|L(0)|+|DL(0)|\leq C_*,
\]
\[
C_*\leq C\left(\|u\|_{L^\infty(B_1)}+
\sup_{0<r<1}r^{1-\alpha}
\left(\frac1{|B_r|}\int_{B_r}|f|^n\right)^{1/n}\right),
\]
where $C$ depends only on $n,\lambda,\Lambda,\alpha$. Thus $u$ is
$C^{1,\alpha}$ at $0$.
\end{theorem}

\section{\Wtwop Estimates}

Assume in this section that $a_{ij},f\in C(B_1)$ and
\[
\lambda|\xi|^2\leq a_{ij}(x)\xi_i\xi_j\leq\Lambda|\xi|^2.
\]

\begin{theorem}\label{hanlin:ch05-thm-5-22}\dcref{Theorem 5.22}\group{ch05-viscosity-solutions}
\source{han-lin-lecture-notes:page-0118,han-lin-lecture-notes:page-0119}
\uses{hanlin:ch05-lem-5-23}
Let $p>n$ and let $u\in C(B_1)$ be a viscosity solution of
\[
a_{ij}D_{ij}u=f\quad\text{in }B_1.
\]
There is $\varepsilon>0$, depending only on $n,\lambda,\Lambda,p$, such
that if
\[
\left(\frac1{|B_r(x_0)|}\int_{B_r(x_0)}
|a_{ij}-a_{ij}(x_0)|^n\right)^{1/n}\leq\varepsilon
\]
for every $B_r(x_0)\subset B_1$, then $u\in W^{2,p}_{\mathrm{loc}}(B_1)$
and
\[
\|u\|_{W^{2,p}(B_{1/2})}
\leq C\left(\|u\|_{L^\infty(B_1)}+\|f\|_{L^p(B_1)}\right),
\]
where $C$ depends only on $n,\lambda,\Lambda,p$.
\end{theorem}
\begin{proof}
\sketch Normalize $u$ and $f$ on balls compactly contained in $B_1$,
apply Lemma~\ref{hanlin:ch05-lem-5-23}, and cover $B_{1/2}$ by finitely
many such balls.
\end{proof}

\begin{lemma}\label{hanlin:ch05-lem-5-23}\dcref{Lemma 5.23}\group{ch05-viscosity-solutions}
\source{han-lin-lecture-notes:page-0119,han-lin-lecture-notes:page-0121,han-lin-lecture-notes:page-0122}
\uses{hanlin:ch05-lem-5-15, hanlin:ch05-lem-5-24}
For each $p>n$ there are $\varepsilon,C>0$, depending only on
$n,\lambda,\Lambda,p$, with the following property. If
$u\in C(B_{8\sqrt n})$ solves
\[
a_{ij}D_{ij}u=f\quad\text{in }B_{8\sqrt n}
\]
in the viscosity sense and
\[
\|u\|_{L^\infty(B_{8\sqrt n})}\leq1,\qquad
\|f\|_{L^p(B_{8\sqrt n})}\leq\varepsilon,
\]
\[
\left(\frac1{|B_r(x_0)|}\int_{B_r(x_0)}
|a_{ij}-a_{ij}(x_0)|^n\right)^{1/n}\leq\varepsilon
\]
for every $B_r(x_0)\subset B_{8\sqrt n}$, then
$u\in W^{2,p}(B_1)$ and $\|u\|_{W^{2,p}(B_1)}\leq C$.
\end{lemma}

\begin{proof}
\sketch Lemma~\ref{hanlin:ch05-lem-5-24} gives a uniform gain in the
measure of contact points. Iterate this gain over dyadic cubes and integrate
the distribution function of the minimal paraboloid opening.
\end{proof}

Equivalently, $x_0\in G_M(u,\Omega)$ when some affine $L$ satisfies,
for all $x\in\Omega$,
\[
L(x)-\frac M2|x-x_0|^2\leq u(x)
\leq L(x)+\frac M2|x-x_0|^2,
\]
with equality at $x_0$. Define
\[
\theta(u,\Omega)(x)=\inf\{M:x\in G_M(u,\Omega)\}\in[0,\infty].
\]
For $p\in(1,\infty]$, if $\theta\in L^p(\Omega)$, then
$D^2u\in L^p(\Omega)$ and
\[
\|D^2u\|_{L^p(\Omega)}\leq2\|\theta\|_{L^p(\Omega)}.
\]
Moreover,
\[
|\{x:\theta(x)>t\}|\leq|A_t(u,\Omega)|\qquad(t>0).
\]

\begin{lemma}\label{hanlin:ch05-lem-5-24}\dcref{Lemma 5.24}\group{ch05-viscosity-solutions}
\source{han-lin-lecture-notes:page-0119,han-lin-lecture-notes:page-0120,han-lin-lecture-notes:page-0121}
\uses{hanlin:ch05-lem-5-15, hanlin:ch05-lem-5-18}
Let $\Omega$ be bounded with $B_{8\sqrt n}\subset\Omega$, and let
$u\in C(\Omega)$ solve
\[
a_{ij}D_{ij}u=f\quad\text{in }B_{8\sqrt n}
\]
in the viscosity sense. For every $\varepsilon_0\in(0,1)$ there are
$M>1$, depending only on $n,\lambda,\Lambda$, and $\varepsilon\in(0,1)$,
depending additionally on $\varepsilon_0$, such that
\[
\|f\|_{L^n(B_{8\sqrt n})}\leq\varepsilon,
\qquad
\|a_{ij}-a_{ij}(0)\|_{L^n(B_{7\sqrt n})}\leq\varepsilon,
\]
\[
G_1(u,\Omega)\cap Q_3\neq\emptyset
\]
imply
\[
|G_M(u,\Omega)\cap Q_1|\geq1-\varepsilon_0.
\]
\end{lemma}

\begin{proof}
\sketch Normalize using the unit supporting paraboloid, compare with the
frozen-coefficient solution from Lemma~\ref{hanlin:ch05-lem-5-18}, and
apply Lemma~\ref{hanlin:ch05-lem-5-15} to the difference.
\end{proof}



\begin{remark}\label{hanlin:ch05-rem-5-25}\dcref{Remark 5.25}\group{ch05-viscosity-solutions}
\source{han-lin-lecture-notes:page-0121}
\uses{hanlin:ch05-lem-5-24}
The conclusion of Lemma~\ref{hanlin:ch05-lem-5-24} remains valid if the
contact hypothesis $G_1(u,\Omega)\cap Q_3\neq\emptyset$ is replaced by
$|u|\leq1$ in $B_{8\sqrt n}$.
\end{remark}
